\documentclass[11pt, a4paper]{article}

% --- PACCHETTI ---
\usepackage[utf8]{inputenc}
\usepackage[italian]{babel}
\usepackage[margin=2.5cm]{geometry} % Imposta i margini della pagina
\usepackage{graphicx}               % Per inserire il logo
\usepackage{tabularx}               % Per tabelle larghe quanto la pagina
\usepackage[table]{xcolor}          % Per colorare le celle delle tabelle
\usepackage[hidelinks]{hyperref}    % Per stile link cliccabili
\usepackage{lmodern}      % Rende i font scalabili all'infinito
\usepackage{anyfontsize}  % Permette di usare dimensioni personalizzate

\begin{document}

% ==============================================================================
% FRONTESPIZIO (PAGINA 1)
% ==============================================================================
\begin{titlepage}
    \centering
    \vspace*{0.5cm}

    % LOGHI
    \includegraphics[width=0.45\textwidth]{../../../../assets/logo_unipd.jpg}
    \hfill
    \includegraphics[width=0.45\textwidth]{../../../../assets/logo.png}\par
    \vspace{0.5cm}

    % NOME DEL GRUPPO
    {\fontsize{45}{50}\sffamily\selectfont \textbf{On-By-One}\par} 
    \vspace{1cm}

    % TITOLO DEL VERBALE
    {\huge \bfseries Verbale interno 2026/10/27 \par}
    
    \vspace{1.5cm}

    % MEMBRI DEL GRUPPO
    {\large
    \textbf{On-By-One}:\par
    \vspace{0.3cm}
    \begin{tabular}{l l}
        Sara Ballardin & 2101048 \\
        Shuhui Chen & 2138001 \\
        Simone Giuliano & 2145591 \\
        Alvise Greggio & 2141004 \\
        Daniele Nuzzaci & 2139039 \\
        Federico Temporin & 2146764 \\
        Tommaso Zorzetto & 2148629 \\
    \end{tabular}
    \par}

    \vfill

    % TABELLA INFORMAZIONI DOCUMENTO E DESTINATARI
    \begin{table}[h]
        \centering
        \arrayrulecolor{gray!50}
        \renewcommand{\arraystretch}{1.5}
        \begin{tabularx}{\textwidth}{|c|c|c|X|}
            \hline
            \rowcolor{gray!10} \multicolumn{4}{|c|}{\textbf{Informazioni documento}} \\
            \hline
            \rowcolor{gray!10} \textbf{Versione} & \textbf{Data} & \textbf{Stato} & \multicolumn{1}{c|}{\textbf{Destinatari}} \\
            \hline
            - & - & Bozza & 
            On-By-One \newline 
            Prof. Tullio Vardanega \newline 
            Prof. Riccardo Cardin \\
            \hline
        \end{tabularx}
        \arrayrulecolor{black}
    \end{table}
    
    \vspace{0.5cm}
    
    % CONTATTI
    {\large Contatti: \href{mailto:onbyone.swe@gmail.com}{onbyone.swe@gmail.com} }
    
\end{titlepage}

\newpage

% ==============================================================================
% REGISTRO DELLE MODIFICHE (PAGINA 2)
% ==============================================================================
\section*{Registro delle modifiche}

\vspace{1cm}
\begin{center}
    \arrayrulecolor{gray!50}
    \renewcommand{\arraystretch}{1.5}
    \begin{tabularx}{\textwidth}{|c|c|c|X|c|}
        \hline
        \rowcolor{gray!10} 
        \multicolumn{1}{|c|}{\textbf{Versione}} & 
        \multicolumn{1}{c|}{\textbf{Data}} & 
        \multicolumn{1}{c|}{\textbf{Autore}} & 
        \multicolumn{1}{c|}{\textbf{Descrizione}} & 
        \multicolumn{1}{c|}{\textbf{Verificatore}} \\
        \hline
        - & 2026/10/28 & Sara Ballardin & Stesura del verbale & - \\
        \hline
        0.1.0 & 2026/10/27 & Sara Ballardin & Bozza iniziale & - \\
        \hline
    \end{tabularx}
    \arrayrulecolor{black}
\end{center}
\vfill
\newpage

% ==============================================================================
% INDICE (PAGINA 3)
% ==============================================================================
\tableofcontents
\newpage

% ==============================================================================
% INFORMAZIONI GENERALI
% ==============================================================================
\section{Ordine del giorno}
\begin{itemize}
    \item Definizione e approvazione della strategia di \textit{branching}.
    \item Definizione e approvazione della policy di \textit{code review}.
    \item Definizione e approvazione delle convenzioni per i \textit{commit}.
    \item Definizione e approvazione delle convenzioni di formattazione.
    \item Regolamentazione della \textit{naming convention} per i file del \textit{repository}.
    \item Definizione del flusso di verifica e approvazione dei documenti.
\end{itemize}

\section{Discussioni}
La riunione è stata indetta con l'obiettivo di definire e formalizzare il \textit{Way of Working} (WoW) del team On-By-One. Si è discusso sulla necessità di stabilire linee guida chiare, condivise e rigorose per la gestione del repository, la stesura dei documenti e le procedure di collaborazione. L'adozione immediata di queste regole garantirà uniformità, tracciabilità ed efficienza nel lavoro del gruppo, semplificando le future fasi di sviluppo e revisione.

\section{Decisioni}
\begin{itemize}
    \item \textbf{VI-2026-10-27.1} Approvata la gestione dei branch di sviluppo: la linea principale (\texttt{main}) è protetta e richiede verifiche automatiche. Le linee di lavoro parallele seguiranno il nome \texttt{<tipo>/<numero-attività>-<descrizione>} (tipi consentiti: \texttt{admin/}, \texttt{docs/}, \texttt{website/}) e verranno eliminate una volta unite al progetto principale tramite \textit{Standard merge}.
    \item \textbf{VI-2026-10-27.2} Approvata la politica di revisione: ogni proposta di modifica (\textit{Pull Request}) deve essere controllata da almeno un membro del gruppo diverso dall'autore. Le correzioni richieste si dividono in \textbf{Bloccanti} (obbligatorie prima di poter unire il lavoro) e \textbf{Non bloccanti} (suggerimenti opzionali o rimandabili).
    \item \textbf{VI-2026-10-27.3} Approvate le regole per i salvataggi (\textit{commit}): i messaggi seguiranno la struttura \texttt{<tipo>: <descrizione>} in italiano, indicando il numero dell'attività di riferimento (es. \texttt{\#24}).
    \item \textbf{VI-2026-10-27.4} Approvate le regole di formattazione: uso del \textbf{grassetto} per i concetti chiave, del \textit{corsivo} solo per la prima comparsa di parole straniere non comuni o termini specialistici, e dello stile \texttt{codice} per nomi di file, comandi e branch. È vietato sottolineare o barrare il testo.
    \item \textbf{VI-2026-10-27.5} Approvate le regole per i nomi dei file: i verbali useranno il formato \texttt{aaaa-mm-dd-vi} per quelli interni e \texttt{aaaa-mm-dd-ve} per quelli esterni. Tutti gli altri documenti avranno nomi in minuscolo, separando le parole con i trattini (es. \texttt{nome-del-file.tex}).
    \item \textbf{VI-2026-10-27.6} Approvato il processo di validazione dei documenti: chi scrive il testo lo lascia in stato di \textbf{Bozza}, senza data o versione definitiva. Chi verifica, se approva il lavoro, inserisce la versione finale e cambia lo stato in \textbf{Approvato}; se invece ci sono problemi, assegna una versione provvisoria, rifiuta la modifica e segnala gli errori da correggere.
\end{itemize}

\section{To-do}
\begin{center}
    \arrayrulecolor{gray!50}
    \renewcommand{\arraystretch}{1.5}
    \begin{tabularx}{\textwidth}{|c|X|c|c|}
        \hline
        \rowcolor{gray!10} 
        \multicolumn{1}{|c|}{\textbf{ID issue}} & 
        \multicolumn{1}{c|}{\textbf{Task}} & 
        \multicolumn{1}{c|}{\textbf{Assegnatario}} & 
        \multicolumn{1}{c|}{\textbf{Scadenza}} \\
        \hline
        \href{https://github.com/On-By-One-Team/On-By-One.github.io/issues/17}{\#17} & Stesura del documento sulla naming convention e sull'approvazione dei documenti & Federico Temporin & 2026/10/28 \\
        \hline
        \href{https://github.com/On-By-One-Team/On-By-One.github.io/issues/21}{\#21} & Aggiornamento del \textit{workflow} per la creazione di PDF & Sara Ballardin & 2026/10/28 \\
        \hline
        \href{https://github.com/On-By-One-Team/On-By-One.github.io/issues/22}{\#22} & Aggiornamento del workflow per il sito web & Daniele Nuzzaci & 2026/10/28 \\
        \hline
    \end{tabularx}
    \arrayrulecolor{black}
\end{center}

\end{document}